% Introducción del informe del proyecto final (modelo original BDL).
\section{Introducción}\label{sec:introduccion}

Retener a un cliente cuesta menos que conseguir uno nuevo, y por eso una \emph{fintech} necesita ordenar a sus clientes
por riesgo de abandono antes de lanzar una campaña de retención de capacidad limitada. Este trabajo aborda ese problema
con COFINFAD \citep{munoz2025}, un conjunto público y sintético de $48.723$ clientes de una \emph{fintech} colombiana con
el detalle de sus transacciones de 2023, cuyo objetivo es un puntaje continuo de propensión al abandono, \cp.

En la versión anterior del proyecto, el modelo propuesto era una \emph{Explainable Boosting Machine} usada desde una
librería. Esta versión la reemplaza por un modelo propio, BDL (\emph{Boosting} con Desplazamiento Latente), que parte de
un hallazgo del análisis exploratorio: el puntaje es un índice suave de los predictores más un escalón constante que
depende de una variable que hay que excluir por fuga de información. BDL modela ese escalón como una clase latente y
aprende el índice con una regla de aprendizaje propia, un EM generalizado cuyo paso M es un paso de \emph{boosting}.

Las contribuciones son tres:
\begin{itemize}
  \item un modelo y su regla de aprendizaje, formulados, derivados y programados desde cero con \texttt{numpy} y
        \texttt{scipy}, con garantías teóricas verificadas por pruebas unitarias (sección~\ref{sec:modelo});
  \item una comparación con todos los modelos del curso y dos del estado del arte, bajo un protocolo registrado antes
        de los experimentos, con las mismas particiones, el mismo preprocesamiento y el mismo presupuesto de ajuste
        (secciones~\ref{sec:referencias} y~\ref{sec:protocolo});
  \item una evaluación estadística con tres pruebas de hipótesis por referencia, tamaños de efecto, cinco conjuntos de
        datos, ablación, sensibilidad, costo e interpretabilidad (sección~\ref{sec:resultados}).
\end{itemize}

La sección~\ref{sec:metodologia} describe los datos y su análisis exploratorio, el pipeline, los modelos de referencia,
el modelo propuesto y el protocolo de evaluación; la sección~\ref{sec:resultados} presenta los resultados, y la
sección~\ref{sec:conclusiones}, las conclusiones.
